%% toc: P15 172 172.7,172.8,172.9,172.10
\section{Sensitivity conclusions under changed dependence and formulation}
\label{sec:annex-p15-w02b-p15-069-assumption-sensitivity}

Every conditioning variable below takes finitely many values or is handled
by product integration. For a discrete conditioning variable, conditional
expectations and variances are the elementary ones on atoms of positive
probability; atoms of probability zero do not affect any expectation.

\begin{annexitem}{annex:p15-w02b-p15-069-coupling}
 {thm:p15-w02b-p15-069-coupling}{Properties of the symmetric binary coupling}
\begin{proof}
Since \(|\rho x_1x_2|\leq1\), every weight \((1+\rho x_1x_2)/4\) is
nonnegative. Summing over the four points, the terms \(\rho x_1x_2\) cancel,
because \(\sum_{x_1,x_2}x_1x_2=0\); the total is \(1\).
For fixed \(x_1\), summing over \(x_2\in\{-1,1\}\) gives
\(P_\rho(X_1=x_1)=1/2\), and by the symmetry of the formula the same holds
for \(X_2\). Hence
\[
 P_\rho(X_2=x_2\mid X_1=x_1)=\frac{1+\rho x_1x_2}{2}.
\]
Its mean is
\(\sum_{x_2}x_2(1+\rho x_1x_2)/2=\rho x_1\), since \(\sum_{x_2}x_2=0\) and
\(x_2^2=1\). As \(X_2^2=1\), the conditional variance is
\(1-\rho^2x_1^2=1-\rho^2\). Both marginals have mean \(0\) and variance
\(1\), and \(E[X_1X_2]=E[X_1E[X_2\mid X_1]]=\rho E[X_1^2]=\rho\), which is
the correlation. Exchanging the coordinates leaves the formula unchanged,
so the same identities hold with the roles exchanged.

Independence means that every weight equals the product \(1/4\) of the
marginals, that is \(\rho x_1x_2=0\) at every point, which happens exactly
when \(\rho=0\). A weight is zero exactly when \(\rho x_1x_2=-1\). For
\(|\rho|<1\) this never happens. For \(\rho=1\) it happens exactly at
\(x_1x_2=-1\), leaving the diagonal; for \(\rho=-1\) exactly at
\(x_1x_2=1\), leaving the antidiagonal.

For \(Q_\rho\), the value \(P(\varepsilon=e)\) equals \((1+\rho e)/2\) for
\(e\in\{-1,1\}\). Since \(a^2=1\), the event \(\{X_1=a,X_3=b\}\) equals
\(\{X_1=a,\varepsilon=ab\}\). Independence of \(X_1\) and \(\varepsilon\)
gives probability \(\tfrac12\cdot\tfrac{1+\rho ab}{2}\), which is the
\(P_\rho\) weight of \((a,b)\). The marginal of \(X_3\) is therefore a fair
sign. Finally \((X_1,X_3)\) is a measurable function of
\((X_1,\varepsilon)\), which is independent of \(X_2\).
\end{proof}
\end{annexitem}

\begin{annexitem}{annex:p15-w02b-p15-069-linear-two}
 {thm:p15-w02b-p15-069-linear-two}{Linear model with coupled inputs}
\begin{proof}
By \cref{thm:p15-w02b-p15-069-coupling}, both inputs have variance \(1\)
and covariance \(\rho\), so
\(V=a_1^2+a_2^2+2a_1a_2\rho\). The same lemma gives
\[
 E[Y\mid X_1]=a_1X_1+a_2\rho X_1=(a_1+a_2\rho)X_1,
\]
whose variance is \((a_1+a_2\rho)^2\), and by exchanging roles
\(\operatorname{Var}(E[Y\mid X_2])=(a_2+a_1\rho)^2\). Given \(X_2\), the
term \(a_2X_2\) is constant, so
\(\operatorname{Var}(Y\mid X_2)=a_1^2\operatorname{Var}(X_1\mid X_2)
 =a_1^2(1-\rho^2)\), a constant. Given \(X_1\), the term \(a_1X_1\) is
constant, so
\(\operatorname{Var}(Y\mid X_1)=a_2^2\operatorname{Var}(X_2\mid X_1)
 =a_2^2(1-\rho^2)\). With two inputs,
\(X_{-1}=X_2\) and \(X_{-2}=X_1\). Dividing by \(V>0\) gives the four
displayed indices.

For item~1,
\[
 (a_1+a_2\rho)^2-(a_2+a_1\rho)^2
 =a_1^2+a_2^2\rho^2-a_2^2-a_1^2\rho^2
 =(a_1^2-a_2^2)(1-\rho^2),
\]
the cross terms \(2a_1a_2\rho\) cancelling, and
\(a_1^2(1-\rho^2)-a_2^2(1-\rho^2)\) is the same expression. For
\(|\rho|<1\) the factor \(1-\rho^2\) is positive, so the sign of each
difference is the sign of \(a_1^2-a_2^2\), which is the sign of the
difference at \(\rho=0\). For \(|\rho|=1\) both differences vanish.

For item~2,
\[
 V(S_1+S_2)=(a_1^2+a_2^2)(1+\rho^2)+4a_1a_2\rho,
\]
and subtracting \(V=a_1^2+a_2^2+2a_1a_2\rho\) leaves
\((a_1^2+a_2^2)\rho^2+2a_1a_2\rho
 =\rho\bigl((a_1^2+a_2^2)\rho+2a_1a_2\bigr)\).
For \(a_1=1\), \(a_2=2\), \(\rho=1/2\), the formulas give \(V=7\),
\(S_1=(1+1)^2/7=4/7\) and \(S_{T_1}=(3/4)/7=3/28\).

For item~3, \(V=5+4\rho\geq1\) for every \(\rho\in[-1,1]\). The values of
\(Y=X_1+2X_2\) at \((1,1)\), \((1,-1)\), \((-1,1)\), \((-1,-1)\) are
\(3,-1,1,-3\). Thus \(Y>2\) exactly at \((1,1)\), which has probability
\((1+\rho)/4\); this is at most \(1/10\) exactly when \(\rho\leq-3/5\).
The bound \(Y\leq1\) fails only at \((1,1)\), so it holds surely exactly
when that point has probability zero, that is, when \(\rho=-1\).
At \(\rho=-3/5\), \(V=13/5\), \((1+2\rho)^2=1/25\), \((2+\rho)^2=49/25\)
and \(1-\rho^2=16/25\), which gives the displayed values
\(1/65\), \(49/65\), \(16/65\), \(64/65\) and their first-order sum
\(50/65=10/13\).
\end{proof}
\end{annexitem}

\begin{annexitem}{annex:p15-w02b-p15-069-duplicate}
 {thm:p15-w02b-p15-069-duplicate}{Duplicated information}
\begin{proof}
By \cref{thm:p15-w02b-p15-069-coupling}, \(X_1\) and \(X_2\) are
independent fair signs, so \(V=4+1=5\) for every \(\rho\).

First-order numerators. \(E[Y\mid X_1]=2X_1+E[X_2]=2X_1\), with variance
\(4\). \(E[Y\mid X_2]=2E[X_1]+X_2=X_2\), with variance \(1\). Since
\((X_1,X_3)\sim P_\rho\) and \(X_2\) is independent of \(X_3\),
\(E[Y\mid X_3]=2E[X_1\mid X_3]+E[X_2]=2\rho X_3\), with variance
\(4\rho^2\).

Total numerators. Each pair of values of \((X_2,X_3)\) has probability
\(1/4\). Because \(X_2\) is independent of \((X_1,X_3)\),
\[
\begin{split}
 P(X_1=a\mid X_2=b,X_3=c)
 &=\frac{P(X_1=a,X_3=c)P(X_2=b)}{P(X_3=c)P(X_2=b)}\\
 &=P(X_1=a\mid X_3=c).
\end{split}
\]
Given \((X_2,X_3)\) the term \(X_2\) is constant, so
\(\operatorname{Var}(Y\mid X_2,X_3)=4\operatorname{Var}(X_1\mid X_3)
 =4(1-\rho^2)\) by the lemma. Given \((X_1,X_3)\), on each atom of
positive probability, \(2X_1\) is constant and \(X_2\) is still a fair
sign, so \(\operatorname{Var}(Y\mid X_1,X_3)=1\). Finally \(Y\) is a
function of \((X_1,X_2)\), so \(\operatorname{Var}(Y\mid X_1,X_2)=0\).
Dividing all six numerators by \(5\) gives the displayed indices.

Item~1 is the case \(\rho=0\). Item~2 compares \(4\rho^2/5\) with
\(1/5\): the first is larger exactly when \(\rho^2>1/4\). Item~3 compares
\(4(1-\rho^2)/5\) with \(1/5\): the second is larger exactly when
\(1-\rho^2<1/4\), that is, \(\rho^2>3/4\); at \(\rho^2=1\),
\(S_{T_1}=0\). Item~4 holds because \(4/5>4(1-\rho^2)/5\) exactly when
\(\rho^2>0\). Item~5 was shown above. At \(\rho=9/10\),
\(4\rho^2/5=81/125\) and \(4(1-\rho^2)/5=4\cdot19/500=19/125\).
\end{proof}
\end{annexitem}

\begin{annexitem}{annex:p15-w02b-p15-069-fixing}
 {thm:p15-w02b-p15-069-fixing}{Substitution error and the total index}
\begin{proof}
Item~1. Let \(G=E[Y\mid X_{-i}]\). If \(h(X_{-i})\) is not square
integrable, then neither is \(Y-h(X_{-i})\), since \(Y\in L^2\), and the
left side is \(+\infty\). Otherwise \(D=G-h(X_{-i})\) is square integrable
and \(X_{-i}\)-measurable, and \(E[Y-G\mid X_{-i}]=0\), so the pull-out
property gives \(E[(Y-G)D]=0\). Hence
\[
 E\bigl[(Y-h(X_{-i}))^2\bigr]=E\bigl[(Y-G)^2\bigr]+E[D^2].
\]
The first term is \(E[\operatorname{Var}(Y\mid X_{-i})]=VS_{T_i}\),
and the second is nonnegative and zero exactly when \(D=0\) almost
surely. The map \(z\mapsto f(c,z)\) is Borel, being the composition of
\(f\) with the Borel map \(z\mapsto(c,z)\), so \(\varphi_i(c)\) is the
case \(h=f(c,\cdot)\).

Item~2. Independence of \(X_i\) and \(X_{-i}\) means that the joint law
is \(P_i\otimes P_{-i}\). By Tonelli,
\[
 \int\varphi_i\,dP_i
 =\int\!\!\int\!\!\int\bigl(f(a,z)-f(c,z)\bigr)^2
   P_i(da)\,P_i(dc)\,P_{-i}(dz).
\]
Write \(F_z=f(\cdot,z)\). Tonelli also gives
\(\int\!\int F_z(a)^2P_i(da)P_{-i}(dz)=E[Y^2]<\infty\), so \(F_z\in
L^2(P_i)\) for \(P_{-i}\)-almost every \(z\). For such \(z\), expanding the
square shows that the inner double integral is
\(2\bigl(\int F_z^2\,dP_i-(\int F_z\,dP_i)^2\bigr)=2v(z)\), where \(v(z)\) is
the variance of \(F_z\) under \(P_i\). Under the product law,
\(m(z)=\int F_z\,dP_i\) and \(\int F_z^2\,dP_i\) are versions of
\(E[Y\mid X_{-i}=z]\) and \(E[Y^2\mid X_{-i}=z]\): for every Borel
\(B\subseteq\mathcal X_{-i}\), Fubini gives
\(E[Y\mathbf 1_B(X_{-i})]=\int_B m\,dP_{-i}\); the same Fubini computation
with the integrable function \(f^2\) in place of \(f\) gives
\(E[Y^2\mathbf 1_B(X_{-i})]=\int_B\bigl(\int F_z^2\,dP_i\bigr)P_{-i}(dz)\).
Since \(\operatorname{Var}(Y\mid X_{-i})=E[Y^2\mid X_{-i}]-E[Y\mid X_{-i}]^2\)
almost surely, \(v(X_{-i})\) is a version of it, and
\(\int v\,dP_{-i}=E[\operatorname{Var}(Y\mid X_{-i})]=VS_{T_i}\). Hence
\(\int\varphi_i\,dP_i=2VS_{T_i}\).

Put \(M=2VS_{T_i}<\infty\). If \(P_i(\varphi_i\leq M)=0\), then
\(\varphi_i-M>0\) almost surely. The sets \(\{\varphi_i-M>1/n\}\) increase
to a set of full measure, so one of them has positive probability
\(p\), and \(\int(\varphi_i-M)\,dP_i\geq p/n>0\). This contradicts
\(\int\varphi_i\,dP_i=M\).

Item~3. Here \(f(c,X_2,X_3)=2c+X_2\), so
\[
 \varphi_1(c)=E\bigl[(2X_1-2c)^2\bigr]
 =4\bigl(E[X_1^2]-2cE[X_1]+c^2\bigr)=4(1+c^2)\geq4,
\]
for every \(\rho\). By \cref{thm:p15-w02b-p15-069-duplicate},
\(VS_{T_1}=5\cdot4(1-\rho^2)/5=4(1-\rho^2)\). If \(\rho^2>1/2\), then
\(2VS_{T_1}=8(1-\rho^2)<4\leq\varphi_1(c)\) for every \(c\). At
\(|\rho|=1\), \(S_{T_1}=0\) while \(\varphi_1(c)\geq4=\tfrac45V\). Since
\(P_1\) is a fair sign, \(\int\varphi_1\,dP_1=4(1+1)=8\), and half of it is
\(4=\tfrac45V\), independently of \(\rho\). At \(\rho=1\) the joint support
has \(x_3=x_1\), while a resampled first coordinate differs from \(X_3\)
with probability \(1/2\); such evaluation points lie outside the support.
\end{proof}
\end{annexitem}

\begin{annexitem}{annex:p15-w02b-p15-069-formulation}
 {thm:p15-w02b-p15-069-formulation}{Curvature family}
\begin{proof}
At \((0,0)\), \(f_\gamma=0\), \(\partial_1f_\gamma=1+2\gamma x_1=1\) and
\(\partial_2f_\gamma=2\).

Item~1. Write \(f_\gamma=g_1(x_1)+g_2(x_2)\) with
\(g_1(x_1)=x_1+\gamma x_1^2\) and \(g_2(x_2)=2x_2\). By independence,
\(E[Y\mid X_1]=g_1(X_1)+E[g_2(X_2)]\), whose variance is
\(\operatorname{Var}(g_1(X_1))\), and
\(\operatorname{Var}(Y\mid X_2)=\operatorname{Var}(g_1(X_1))\) because
the conditional law of \(X_1\) given \(X_2\) is its marginal. Thus
\(S_1=S_{T_1}\), and likewise \(S_2=S_{T_2}=\operatorname{Var}(g_2(X_2))/V\).
For \(X\) uniform on \([-1,1]\),
\(E[X]=E[X^3]=0\), \(E[X^2]=1/3\) and \(E[X^4]=1/5\). Hence
\[
\begin{split}
 \operatorname{Var}(g_1(X_1))
 &=E\Bigl[\bigl(X_1+\gamma(X_1^2-\tfrac13)\bigr)^2\Bigr]\\
 &=\frac13+2\gamma\Bigl(E[X_1^3]-\frac{E[X_1]}3\Bigr)
  +\gamma^2\Bigl(\frac15-\frac19\Bigr)
 =\frac{15+4\gamma^2}{45},
\end{split}
\]
and \(\operatorname{Var}(g_2(X_2))=4/3=60/45\). By independence
\(V=(75+4\gamma^2)/45>0\), which gives both ratios. Then \(S_2>S_1\) exactly
when \(60>15+4\gamma^2\), that is \(\gamma^2<45/4\), and
\(\sqrt{45/4}=3\sqrt5/2\).

Item~2. If \(\gamma\leq0\), then \(\gamma x_1^2\leq0\) and
\(f_\gamma\leq x_1+2x_2\leq3\) on the square. If \(\gamma>0\), then
\(f_\gamma(1,1)=3+\gamma>3\). By continuity there is \(\delta\in(0,1)\) with
\(f_\gamma>3\) on \([1-\delta,1]^2\), a set of probability \(\delta^2/4>0\)
under the uniform law of density \(1/4\).

Item~3. With \(4\gamma^2=64\), item~1 gives \(79/139\) and \(60/139\).
On \([-1,1]\), \(x_1+4x_1^2\leq|x_1|+4x_1^2\leq5\), with equality at
\(x_1=1\), and \(2x_2\leq2\), with equality at \(x_2=1\); so the maximum
is \(7\). The function \(t\mapsto t+4t^2\) is increasing for \(t>-1/8\).
On \((3/4,1]\times(0,1]\), therefore,
\(x_1+4x_1^2>3/4+4\cdot9/16=3\) and \(2x_2>0\), so \(f_4>3\). This
rectangle has probability \(\tfrac18\cdot\tfrac12=\tfrac1{16}\).
\end{proof}
\end{annexitem}
